\documentclass{article}
\usepackage[UTF8]{ctex}
\usepackage{amsmath,amssymb,amsfonts}
\usepackage{enumitem}
\usepackage{geometry}
\geometry{a4paper, margin=1in}
% === 错题本排版优化（不参与题目标识） ===
\setlength{\parindent}{0pt}
\setlength{\parskip}{0.35em plus 0.10em minus 0.08em}
\setlength{\emergencystretch}{3em}
\setlist{itemsep=0.25em, topsep=0.40em, parsep=0pt}
\allowdisplaybreaks[2]
% === 错题本排版优化结束 ===
\begin{document}
\section*{西安电子科技大学2024年高等代数真题}
\section*{180．西安电子科技大学2024年高等代数试卷}

\subsection{题目1}
一．设矩阵 \(A \in \mathbb{C}^{n \times n}\) ，复系数多项式 \(f(x)=h(x) g(x)\) ，其中 \((h(x), g(x))=1\) ，证明：\(f(A)=O\)的充要条件是 \(\operatorname{rank}(h(A))+\operatorname{rank}(g(A))=n\) 。

\subsection{题目2}
二．设 \(x_{1}, x_{2}, \cdots, x_{n}\) 是互不相同的 \(n\) 个数．\\
（1）记 \(|A(t)|=\left|\begin{array}{cccc}x_{1}+t & x_{1}^{2}+t & \cdots & x_{1}^{n}+t \\ x_{2}+t & x_{2}^{2}+t & \cdots & x_{2}^{n}+t \\ \vdots & \vdots & & \vdots \\ x_{n}+t & x_{n}^{2}+t & \cdots & x_{n}^{n}+t\end{array}\right|\) ，其中 \(t\) 是参数。证明：

\[
|A(t)|=|A(0)|+t \sum_{i, j=1}^{n} A_{i j}
\]

其中 \(A_{i j}\) 是 \(x_{i}^{j}\) 在 \(|A(0)|\) 中的代数余子式．\\
（2）计算

\[
\left|\begin{array}{cccc}
x_{1}+1 & x_{1}^{2}+1 & \cdots & x_{1}^{n}+1 \\
x_{2}+1 & x_{2}^{2}+1 & \cdots & x_{2}^{n}+1 \\
\vdots & \vdots & & \vdots \\
x_{n}+1 & x_{n}^{2}+1 & \cdots & x_{n}^{n}+1
\end{array}\right|
\]

\subsection{题目3}
三。设 \(A\) 是 \(n\) 阶矩阵，\(b \in \mathbb{R}^{n}\) ，解答如下问题：

\begin{itemize}
  \item[（1）] 若 \(A\) 为实矩阵，证明：线性方程组 \(A^{\mathrm{T}} A X=A^{\mathrm{T}} b\) 有解。
  \item[（2）] 若 \(A\) 为复矩阵，线性方程组 \(A^{\mathrm{T}} A X=A^{\mathrm{T}} b\) 是否一定有解？为什么？
\end{itemize}

（3）若 \(A\) 为复矩阵，线性方程组 \(A^{H} A X=A^{H} b\) 是否一定有解？为什么？

\subsection{题目4}
四．求矩阵

\[
H=\left(\begin{array}{ccccc}
1 & -b & & & \\
& 1 & -b & & \\
& & \ddots & & \\
& & & 1 & -b \\
& & & & 1
\end{array}\right)
\]

的逆．

\subsection{题目5}
五 ．解答如下问题：\\
（1）证明：实二次型 \(f\left(x_{1}, x_{2}, \cdots, x_{n}\right)=X^{\mathrm{T}} A X\) 的秩等于矩阵 \(A\) 非零特征值的个数，其中 \(A\) 为实对称矩阵，\(X=\left(x_{1}, x_{2}, \cdots, x_{n}\right)^{\mathrm{T}}\) 。\\
（2）化二次型 \(\sum_{i=1}^{n}\left(x_{i}-\bar{x}\right)^{2}\) 为标准形，其中 \(\bar{x}=\frac{x_{1}+x_{2}+\cdots+x_{n}}{n}\) ．

\subsection{题目6}
六。已知 \(A, B\) 分别为 \(m \times n, n \times s\) 矩阵，证明： \(\operatorname{rank}(A B)=\operatorname{rank}(A)+\operatorname{rank}(B)-n\) 的充要条件是\(N(A) \subseteq R(B)\) ，其中 \(N(A), R(B)\) 分别为矩阵 \(A\) 的零空间和矩阵 \(B\) 的值域空间．

\subsection{题目7}
七．设 \(n\) 阶复矩阵 \(A, B\) 均可相似对角化，且 \(A B=B A\) ，证明：存在可逆矩阵 \(P\) ，使得 \(P^{-1} A P, P^{-1} B P\)都是对角矩阵。

\subsection{题目8}
八．证明：\(n\) 阶矩阵 \(A\) 是幂等矩阵的充要条件是 \(\operatorname{rank}(A)+\operatorname{rank}\left(I_{n}-A\right)=n\) 。

\subsection{题目9}
九．设 \(A\) 为 \(n\) 阶方阵， 1 不是 \(A\) 的不变因子， 3 是 \(A\) 的一个特征值，证明：\(A=3 I_{n}\) ．

\subsection{题目10}
十．设 \(V\) 是由 \(\mathbb{R}^{3 \times 3}\) 中所有斜对称矩阵组成的线性空间，对任意的 \(A, B \in V\) ，定义 \((A, B)=\) \(\frac{1}{2} \operatorname{tr}\left(A B^{\mathrm{T}}\right)\) ．\\
（1）证明：（，）是 \(V\) 上的一个内积．\\
（2）令 \(\sigma: \mathbb{R}^{3} \rightarrow V\) 如下：

\[
\sigma:\left(\begin{array}{l}
x_{1} \\
x_{2} \\
x_{3}
\end{array}\right) \mapsto\left(\begin{array}{ccc}
0 & x_{1} & x_{2} \\
-x_{1} & 0 & x_{3} \\
-x_{2} & -x_{3} & 0
\end{array}\right) .
\]

证明：\(\sigma\) 是欧几里得空间 \(\mathbb{R}^{3}\)（指定标准内积）到 \(V\) 的一个同构映射，并求 \(V\) 的一组标准正交基．

\subsection{题目11}
十一．设 \(W\) 是向量组 \(\alpha_{1}, \alpha_{2}, \alpha_{3}\) 张成的线性子空间，其中 \(\alpha_{1}=(1,0,2,1)^{\mathrm{T}}, \alpha_{2}=(2,1,2,3)^{\mathrm{T}}, \alpha_{3}=\) \((0,1,-2,1)^{\mathrm{T}}\) ，求 \(W\) 正交补 \(W^{\perp}\) 的标准正交基。
\end{document}
